\chapter{64-Mbit SDRAM Memory}
\label{ch:sdram}

\section{Why an FPGA Needs External Memory}

Flip-flops and BSRAM block memory are very fast, but there is little of them.
For an image, sound, or a large array of data, you need a larger memory. This
is exactly why the GW2AR-18 package includes an SDRAM chip of
\textbf{64 Mbit = 8 MB}.

\begin{figure}[H]
\centering
\begin{tikzpicture}
\begin{axis}[
  xbar, bar width=5mm, width=0.8\textwidth, height=5.8cm,
  xmode=log, log origin=infty, xmin=1e3, xmax=3e8,
  xlabel={Capacity, bits (logarithmic scale)},
  symbolic y coords={Flip-flops, SSRAM, BSRAM, SDRAM},
  ytick=data, y tick label style={font=\sffamily\small},
  nodes near coords, nodes near coords style={font=\sffamily\small, anchor=west},
  point meta=explicit symbolic,
  label style={font=\sffamily\small}, x tick label style={font=\small},
  axis lines*=left, grid=major, grid style={FpgaGray!20}, enlarge y limits=0.2,
]
  \addplot[fill=FpgaOrange!70, draw=FpgaOrange] coordinates {
    (15552,Flip-flops) [$\approx$1.9 KB] (41472,SSRAM) [$\approx$5 KB]
    (847872,BSRAM) [$\approx$103 KB] (67108864,SDRAM) [8 MB]};
\end{axis}
\end{tikzpicture}
\caption{How much data can be stored in the different kinds of memory of the Tang Nano 20K}
\end{figure}

For example, one image frame of $640\times480$ pixels at 16 bits per pixel
takes
\[
  640 \cdot 480 \cdot 16 = 4\,915\,200 \text{ bits} \approx 600 \text{ KB},
\]
which is six times more than all the FPGA's block memory, but fits easily in
SDRAM: more than ten such frames fit there.

\begin{diybox}[How much data fits in internal memory]
\textbf{Task.} Estimate whether your data array fits in block memory, and see how the synthesizer places a small memory.

\textbf{In Logra.} The From an example card: RAM: write and read, and Eight-word ROM. Code tab $\to$ Synthesize; the Gowin tab or the cell list in the report.

\textbf{Expected result.} The array \code{reg [7:0] mem [0:15]} is placed by the synthesizer as a memory (not a hundred flip-flops), while for a $640\times 480\times 16$-bit frame ($\approx 4.9$ Mbit) the block memory (about 0.85 Mbit) is not enough.

\textbf{How to check.} Make a table ``array $\to$ size in bits $\to$ fits in BSRAM or not'' and compare it with the memory-size figure at the beginning of the chapter.

\stars{1} repeat: the size of an 8-word by 8-bit ROM. \stars{2} modify: make the array 256$\times$8 and see how the report changed. \stars{3} project: Security Panel stage---event log; acceptance criterion: you have calculated how many 8-bit events fit in internal memory and how many fit in SDRAM.
\end{diybox}

\section{How DRAM Works}

A dynamic memory (DRAM) cell consists of just \textbf{one transistor and one
capacitor}. A charged capacitor stores 1, a discharged one stores 0. That is
why DRAM is so dense and cheap. But it has two inconveniences:
\begin{itemize}
  \item the capacitor gradually \textbf{discharges}, so every row
        must be periodically rewritten (\term{refresh})---all
        rows within 64~ms;
  \item reading destroys the charge, so a row is first copied into a
        fast buffer (this is called \term{opening} it), and when the work is done it is written
        back (\term{closing} it).
\end{itemize}

\begin{figure}[H]
\centering
\begin{circuitikz}[scale=0.9, transform shape]
  \draw (3,1.5) node[nmos] (T) {};
  \draw (0,0 |- T.gate) node[left, font=\sffamily\small] {row (WL)} -- (T.gate);
  \draw (T.drain) -- (T.drain |- 0,3.4) node[above, font=\sffamily\small] {column (BL)};
  \draw (T.source) to[C, l_=$C$] ++(0,-1.6) node[ground]{};
  \node[font=\sffamily\small, align=left, anchor=west] at (4.6,1.2) {1 bit = 1 transistor + 1 capacitor\\[1mm]
  \textcolor{FpgaRed}{charge leaks away} $\Rightarrow$ refresh is needed};
\end{circuitikz}
\caption{A DRAM cell (simplified)}
\end{figure}

\begin{diybox}[Refresh counter]
\textbf{Task.} All 2048 rows must be refreshed within 64~ms. Build a module that issues a \code{refresh\_req} request with the right period.

\textbf{In Logra.} Code tab: a counter with the constant $N = 100\,\text{MHz} \cdot 31.25\,\mu\text{s} = 3125$ (for the simulation, use $N=25$). Waveforms: \code{cnt}, \code{refresh\_req}.

\textbf{Expected result.} One request per 3125 clock cycles (per 25 in the simulation); \code{refresh\_req} lasts one clock cycle.

\textbf{How to check.} Waveforms: there is exactly one pulse per period $N$; $2048\cdot 31.25$~$\mu$s $=64$~ms.

\stars{1} repeat: a period of 31.25 $\mu$s. \stars{2} modify: the SDRAM frequency is 133~MHz---recalculate $N$. \stars{3} project: Security Panel stage---log refresh; acceptance criterion: a refresh request is never skipped, even while the log is being written.
\end{diybox}

\section{Memory Organization}

The SDRAM in the GW2AR-18 has a 32-bit data bus and is organized as
$4 \text{ banks} \times 2048 \text{ rows} \times 256 \text{ columns}
\times 32 \text{ bits}$:
\[
  4 \cdot 2048 \cdot 256 \cdot 32 = 67\,108\,864 \text{ bits} = 64 \text{ Mbit}.
\]

\begin{figure}[H]
\centering
\begin{tikzpicture}[x=1cm, y=1cm]
  \foreach \b/\c in {3/FpgaBlue!15, 2/FpgaBlue!25, 1/FpgaBlue!35, 0/FpgaBlue!50} {
    \begin{scope}[shift={(\b*0.45,\b*0.45)}]
      \fill[\c, draw=FpgaBlue] (0,0) rectangle (5,3.2);
      \ifnum\b>0 \node[font=\sffamily\scriptsize, anchor=north east] at (5,3.2) {bank \b};\fi
    \end{scope}
  }
  \foreach \y in {0.4,0.8,...,3.0} \draw[FpgaBlue!40] (0,\y) -- (5,\y);
  \foreach \x in {0.5,1.0,...,4.6} \draw[FpgaBlue!40] (\x,0) -- (\x,3.2);
  \node[font=\sffamily\scriptsize, anchor=north east, fill=FpgaBlue!50, inner sep=1pt] at (5,3.2) {bank 0};
  \fill[FpgaOrange!80] (0,1.6) rectangle (5,2.0);
  \fill[FpgaRed] (2.5,1.6) rectangle (3.0,2.0);
  \draw[decorate, decoration={brace, amplitude=5pt}, thick] (-0.2,0) -- (-0.2,3.2) node[midway, left=6pt, font=\sffamily\small, align=right] {2048 rows\\\scriptsize address A[10:0]};
  \draw[decorate, decoration={brace, amplitude=5pt, mirror}, thick] (0,-0.2) -- (5,-0.2) node[midway, below=6pt, font=\sffamily\small, align=center] {256 columns, address A[7:0]};
  \node[font=\sffamily\small, align=left, anchor=west] at (7.4,2.6) {\textcolor{FpgaOrange}{\rule{3mm}{3mm}} open row};
  \node[font=\sffamily\small, align=left, anchor=west] at (7.4,1.9) {\textcolor{FpgaRed}{\rule{3mm}{3mm}} one 32-bit word};
  \node[font=\sffamily\small, align=left, anchor=west] at (7.4,1.2) {the bank is selected by BA[1:0]};
\end{tikzpicture}
\caption{SDRAM organization: bank, row, column}
\end{figure}

\begin{diybox}[Split an address into bank, row, and column]
\textbf{Task.} A word is addressed by 21 bits. Split the address into a bank (2 bits), a row (11 bits), and a column (8 bits).

\textbf{In Logra.} Code tab: a module with \code{input [20:0] addr} and \code{output [1:0] ba, output [10:0] row, output [7:0] col}; the line \code{assign \{ba, row, col\} = addr;}.

\textbf{Expected result.} For \code{addr = 0x1FFFFF} you get \code{ba = 3}, \code{row = 0x7FF}, \code{col = 0xFF}; for \code{addr = 0x000100}---\code{ba = 0}, \code{row = 0}, \code{col = 0x00} and \code{row = 1} (the next row).

\textbf{How to check.} Waveforms: apply several addresses and compare the fields; on the RTL tab there is not a single gate---only wires, because this is just a renaming of bits.

\stars{1} repeat: the split for \code{0x1FFFFF}. \stars{2} modify: change the field order to \code{\{row, ba, col\}} and explain why interleaving banks speeds up the exchange. \stars{3} project: Security Panel stage---log address; acceptance criterion: the write-address counter reaches $2^{21}-1$ and returns to 0.
\end{diybox}

\section{SDRAM Commands}

Unlike BSRAM, you cannot simply apply an address to SDRAM and get data back.
The chip is controlled by \term{commands}, which are encoded by a combination
of the signals \code{CS\#}, \code{RAS\#}, \code{CAS\#}, \code{WE\#} (the hash
sign means active low) on the clock edge.

\begin{table}[H]
\centering
\caption{The main SDRAM commands}
\begin{tabular}{l c c c c l}
\toprule
\textbf{Command} & \textbf{CS\#} & \textbf{RAS\#} & \textbf{CAS\#} & \textbf{WE\#} & \textbf{Action} \\
\midrule
NOP        & 0 & 1 & 1 & 1 & do nothing \\
ACTIVE     & 0 & 0 & 1 & 1 & open a row in a bank \\
READ       & 0 & 1 & 0 & 1 & read from the open row \\
WRITE      & 0 & 1 & 0 & 0 & write to the open row \\
PRECHARGE  & 0 & 0 & 1 & 0 & close a row \\
AUTO REFRESH & 0 & 0 & 0 & 1 & refresh \\
LOAD MODE  & 0 & 0 & 0 & 0 & configuration (CAS latency, burst length) \\
\bottomrule
\end{tabular}
\end{table}

\begin{figure}[H]
\centering
\begin{tikztimingtable}[timing/slope=0.1, xscale=1.9, yscale=1.4, timing/font=\sffamily\small, timing/rowdist=3.3ex]
  CLK      & 9{1L 1H} \\
  Command  & 2D{ACTIVE} 2D{NOP} 2D{READ} 2D{NOP} 2D{NOP} 2D{PRE} 6D{NOP} \\
  Address  & 2D{row} 2U 2D{col.} 12U \\
  DQ[31:0] & 10Z 2D{D0} 2D{D1} 4Z \\
\extracode
  \draw[FpgaOrange, <->, thick] (1,-6.6) -- (5,-6.6) node[midway, below, font=\sffamily\scriptsize] {$t_\text{RCD}$};
  \draw[FpgaPurple, <->, thick] (5,-6.6) -- (9,-6.6) node[midway, below, font=\sffamily\scriptsize] {CAS latency = 2};
\end{tikztimingtable}
\caption{Reading from SDRAM: open a row (ACTIVE), wait $t_\text{RCD}$,
issue the READ command, and receive the data after CAS latency clock cycles}
\label{fig:sdram_read}
\end{figure}

\begin{diybox}[SDRAM command encoder]
\textbf{Task.} Build a module that, given a command number, produces the four signals \code{cs\_n}, \code{ras\_n}, \code{cas\_n}, \code{we\_n} according to the command table.

\textbf{In Logra.} Code tab: \code{case (cmd)} with branches NOP, ACTIVE, READ, WRITE, PRECHARGE, REFRESH, LOAD MODE. Waveforms: step through \code{cmd = 0\ldots6}.

\textbf{Expected result.} NOP $=0111$, ACTIVE $=0011$, READ $=0101$, WRITE $=0100$, PRECHARGE $=0010$, AUTO REFRESH $=0001$, LOAD MODE $=0000$ (the order is \code{cs\_n, ras\_n, cas\_n, we\_n}).

\textbf{How to check.} Waveforms: compare all seven combinations with the table; the \code{default} branch must produce NOP.

\stars{1} repeat: ACTIVE and READ. \stars{2} find the bug: remove the \code{default} branch and explain where the latch on the RTL came from. \stars{3} project: Security Panel stage---log commands; acceptance criterion: all seven commands are checked on the waveform with no gaps.
\end{diybox}

\section{The SDRAM Controller}

Remembering all the commands and delays for every memory access is
inconvenient. That is why a \term{controller} is placed between the user logic
and the SDRAM---a module (usually a finite-state machine, as in project 5) that:
\begin{itemize}
  \item initializes the chip after power-up;
  \item turns simple requests like ``read the word at address X'' into a
        sequence of commands ACTIVE $\to$ READ $\to$ PRECHARGE;
  \item performs the refresh on time by itself.
\end{itemize}

\begin{figure}[H]
\centering
\begin{tikzpicture}[node distance=24mm]
  \node[gblock, minimum height=24mm, minimum width=24mm] (user) {Your logic\\\scriptsize (video, sound\ldots)};
  \node[hblock, right=of user, minimum height=24mm, minimum width=26mm] (ctrl) {SDRAM\\controller\\\scriptsize state machine};
  \node[oblock, right=of ctrl, minimum height=24mm, minimum width=20mm] (mem) {SDRAM\\64 Mbit};
  \draw[arr] ([yshift=6mm]user.east) -- node[above, lbl] {addr, wr, rd} ([yshift=6mm]ctrl.west);
  \draw[arr] ([yshift=0mm]user.east) -- node[above, lbl] {din[31:0]} ([yshift=0mm]ctrl.west);
  \draw[arr] ([yshift=-6mm]ctrl.west) -- node[below, lbl] {dout, busy} ([yshift=-6mm]user.east);
  \draw[arr] ([yshift=6mm]ctrl.east) -- node[above, lbl] {commands} ([yshift=6mm]mem.west);
  \draw[arr] ([yshift=0mm]ctrl.east) -- node[above, lbl] {A, BA} ([yshift=0mm]mem.west);
  \draw[arr, <->] ([yshift=-6mm]ctrl.east) -- node[below, lbl] {DQ[31:0]} ([yshift=-6mm]mem.west);
  \node[lbl, below=2mm of user] {simple requests};
  \node[lbl, below=2mm of mem] {complex protocol};
\end{tikzpicture}
\caption{The controller hides the complexity of the SDRAM protocol}
\end{figure}

\subsection{Ports of the Built-in Memory}

In a Gowin project, the built-in SDRAM is connected to the top-level module
through ports with special (reserved) names. You do not need to specify pin
numbers in the \code{.cst} file for them: the connection is inside the
package.

\begin{vcode}
module top (
    input  wire        clk,
    // --- Built-in SDRAM (do not change the names!) ---
    output wire        O_sdram_clk,
    output wire        O_sdram_cke,
    output wire        O_sdram_cs_n,    // CS#
    output wire        O_sdram_ras_n,   // RAS#
    output wire        O_sdram_cas_n,   // CAS#
    output wire        O_sdram_wen_n,   // WE#
    output wire [3:0]  O_sdram_dqm,     // byte mask
    output wire [10:0] O_sdram_addr,    // row/column address
    output wire [1:0]  O_sdram_ba,      // bank number
    inout  wire [31:0] IO_sdram_dq      // data
);
\end{vcode}

\begin{infobox}[Where to get a controller]
\begin{itemize}
  \item Gowin EDA has a ready-made IP block: \textbf{Tools $\to$ IP Core Generator
        $\to$ Memory Interface $\to$ SDRAM Controller}. The wizard
        generates a module with a simple interface.
  \item There are open-source controllers written for the Tang Nano 20K
        by enthusiasts; they are used, for example, in the retro-console projects NESTang and
        SNESTang.
  \item A good advanced exercise is to write your own controller,
        following Fig.~\ref{fig:sdram_read} and the chip documentation.
\end{itemize}
\end{infobox}

\begin{diybox}[Controller skeleton: ACTIVE, READ, PRECHARGE]
\textbf{Task.} Describe a state machine that, on a \code{rd} request, issues the sequence from Fig.~\ref{fig:sdram_read}: ACTIVE, wait $t_\text{RCD}=2$, READ, wait CAS latency $=2$, PRECHARGE.

\textbf{In Logra.} Code tab: the states IDLE, ACT, WAIT1, RD, WAIT2, PRE (modeled on the traffic-light state machine from project 5). The output is the command number from the previous task. Waveforms: the \code{rd} pulse.

\textbf{Expected result.} After \code{rd}, the commands go in the order ACTIVE, NOP, READ, NOP, NOP, PRECHARGE, and then back to IDLE.

\textbf{How to check.} Waveforms: put the cursor on READ and make sure there are exactly two clock cycles between ACTIVE and READ. RTL: a state register and transition logic.

\stars{1} repeat: the command sequence from the figure. \stars{2} modify: make $t_\text{RCD}=3$ a \code{parameter}. \stars{3} project: Security Panel stage---log controller; acceptance criterion: writing one word to the log takes the sequence ACTIVE, WRITE, PRECHARGE without violating the delays.
\end{diybox}

\section{Bandwidth}

At an SDRAM clock frequency of, say, 100~MHz and a 32-bit bus, the peak
transfer rate is
\[
  B = 100 \cdot 10^6 \ \text{Hz} \cdot 32 \ \text{bits} = 3.2 \ \text{Gbit/s} = 400 \ \text{MB/s}.
\]
The real rate is lower because of opening and closing rows and because of
refresh. The longer the bursts of consecutive addresses, the closer the rate
gets to the peak.

\begin{figure}[H]
\centering
\begin{tikzpicture}
\begin{axis}[
  width=0.82\textwidth, height=5.6cm,
  xlabel={Burst length, words}, ylabel={Effective rate, MB/s},
  xmin=1, xmax=256, xmode=log, log basis x=2, ymin=0, ymax=420,
  xtick={1,2,4,8,16,32,64,128,256}, xticklabels={1,2,4,8,16,32,64,128,256},
  label style={font=\sffamily\small}, tick label style={font=\sffamily\small},
  grid=major, grid style={FpgaGray!20}, axis lines*=left,
]
  \addplot[very thick, FpgaBlue, samples at={1,2,4,8,16,32,64,128,256}, mark=*] {400*x/(x+6)};
  \addplot[thick, dashed, FpgaRed, domain=1:256] {400};
  \node[font=\sffamily\small, text=FpgaRed, anchor=north west] at (axis cs:1.3,388) {peak 400 MB/s};
\end{axis}
\end{tikzpicture}
\caption{An estimate of the transfer rate at 100 MHz: each burst has about
6 ``overhead'' clock cycles (ACTIVE, $t_\text{RCD}$, CAS latency, PRECHARGE)}
\end{figure}

\begin{ideabox}
SDRAM gives a lot of memory (8 MB), but it needs a controller: rows must be
opened and closed, and the memory must be refreshed regularly. The highest
speed is obtained when data is read and written in long consecutive bursts.
\end{ideabox}

\begin{diybox}[Measure the effective rate]
\textbf{Task.} Calculate the effective transfer rate for bursts of 1, 8, 64, and 256 words.

\textbf{In Logra.} Waveforms tab, on the model from the previous task: add a clock-cycle counter that runs while \code{busy = 1}, and compare with the formula $B(L)=400\cdot L/(L+6)$.

\textbf{Expected result.} For $L=1$ you get about 57~MB/s, for 8---229, for 64---366, for 256---391~MB/s: the longer the burst, the closer the rate to the peak of 400.

\textbf{How to check.} Substitute the values into the formula and plot the points on the chapter's graph.

\stars{1} repeat: $L=64$. \stars{2} modify: take 8 overhead clock cycles instead of 6 and recalculate. \stars{3} project: Security Panel stage---log capacity and speed; acceptance criterion: one formula estimates how long it takes to write a log of 100\,000 events.
\end{diybox}

\section{Questions and Problems}
\label{sec:tasks-sdram}

The answers are collected in Appendix~\ref{app:answers}.

\begin{taskbox}[Chapter problems]
\begin{enumerate}
  \item \stars{1}~How many 32-bit words does the SDRAM from this chapter store? How many bits must a word address have?
  \item \stars{1}~What is the peak rate at a frequency of 133~MHz and a 32-bit bus?
  \item \stars{2}~All 2048 rows must be refreshed within 64~ms. What is the period of the AUTO REFRESH commands? How many clock cycles is that at 100~MHz?
  \item \stars{2}~Find the bug. A controller issues READ on the clock cycle right after ACTIVE, although $t_\text{RCD}=2$.
        What will happen? How can the state machine be fixed?
  \item \stars{3}~Project. Write a module that splits a 21-bit word address into a bank, a row, and a column
        and encodes the SDRAM commands. Acceptance criterion: on the waveform, all seven commands and several addresses give
        values from the tables of the chapter.
  \item \stars{3}~Security Panel project, stage ``log.'' Describe how the event log will be stored in external
        memory. Acceptance criterion: the width of the address counter, the refresh period, and an estimate of the log fill time are given.
\end{enumerate}
\end{taskbox}
